\documentclass[11pt]{article}

\usepackage[margin=1in]{geometry}
\usepackage{amsmath,amssymb,amsthm}
\usepackage{graphicx}
\usepackage[hidelinks]{hyperref}

\newtheorem{theorem}{Theorem}

\newcommand{\QSL}{\mathit{QSL}}
\newcommand{\Rep}{\operatorname{Rep}}
\newcommand{\Cyc}{\operatorname{Cyc}}
\newcommand{\card}{\operatorname{card}}
\newcommand{\ellp}{\ell_p^2}

\title{A Counterexample to the Converse Between Two \(p\)-Universal Representation Notions}
\author{arXiv automatic math run \texttt{fa\_banach\_001}}
\date{June 19, 2026}

\begin{document}
\maketitle

\section*{Status}

This packet gives a candidate counterexample to the first open remark after
Definition 3.7 in Neufang--Runde \cite{NR2007}.  It is a counterexample to a
converse between two definitions of \(p\)-universal representation.  It does
not address the later open questions in the source paper about Herz--Schur
multipliers or complete isometries.

\section*{Source Question}

The source paper defines a representation \((\pi_u,E_u)\in\Rep_p(G)\) to be
\(p\)-universal if it contains every small representation in \(\Rep_p(G)\).
Immediately afterwards, the authors write:

\begin{center}
\includegraphics[width=\textwidth]{figures/open_problem_crop.png}
\end{center}

Thus the extracted question is whether the older \(p\)-universal notion from
Runde \cite[Definition 4.5]{Runde2004}, where one contains every cyclic
representation, implies the stronger small-representation universality of
Definition 3.7 in \cite{NR2007}.

\section*{Definitions Used}

For a locally compact group \(G\), \(\Rep_p(G)\) denotes the representations
of \(G\) by invertible isometries on \(\QSL_p\)-spaces.  A representation
\((\rho,F)\) is contained in \((\pi,E)\) if it is isometrically equivalent to
a subrepresentation.

In Runde's older sense, \((\pi,E)\) is \(p\)-universal if it contains every
cyclic representation in \(\Rep_p(G)\), where cyclic means that some vector
\(\xi\) has dense cyclic span \(\overline{\pi(L_1(G))\xi}=E\)
\cite[Definition 4.5]{Runde2004}.  In Neufang--Runde's stronger sense,
\((\pi,E)\) is \(p\)-universal if it contains every small representation,
where small means cardinality at most
\(\card(L_1(G))^{\aleph_0}\) \cite[Definition 3.7]{NR2007}.

\section*{Proof Intuition}

For the trivial group \(G=\{e\}\), the algebra \(L_1(G)\) is just
\(\mathbb C\), and its action in any representation is scalar
multiplication.  Therefore a cyclic representation can only be
one-dimensional.  The one-dimensional trivial representation consequently
contains all cyclic representations.  But the stronger 2007 definition asks
for every small representation, and the trivial representation on
\(\ell_p^2\) is small.  A one-dimensional space cannot contain it.

\begin{theorem}
For every \(p\in(1,\infty)\), the converse in the first remark after
Definition 3.7 of \cite{NR2007} is false.  More precisely, there is a
representation that is \(p\)-universal in the sense of
\cite[Definition 4.5]{Runde2004} but is not \(p\)-universal in the sense of
\cite[Definition 3.7]{NR2007}.
\end{theorem}

\begin{proof}
Let \(G=\{e\}\) be the trivial locally compact group.  Then
\(L_1(G)\cong\mathbb C\).  Every representation of \(G\) on a Banach space is
the trivial representation, and the integrated \(L_1(G)\)-action is scalar
multiplication:
\[
  \pi(a)\xi = a\xi \qquad (a\in\mathbb C,\ \xi\in E).
\]

Let \((\pi,E)\in\Rep_p(G)\) be cyclic in Runde's sense.  Then there is
\(\xi\in E\) such that
\[
  E=\overline{\pi(L_1(G))\xi}=\overline{\mathbb C\xi}.
\]
Hence \(E\) is either zero-dimensional or one-dimensional.  Every nonzero
one-dimensional complex Banach space is isometrically isomorphic to
\(\mathbb C\), and the group action is necessarily trivial.  Therefore the
one-dimensional trivial representation
\[
  (\mathbf 1,\mathbb C)\in\Rep_p(G)
\]
contains every cyclic representation in \(\Rep_p(G)\).  It is
\(p\)-universal in the sense of \cite[Definition 4.5]{Runde2004}.

We now check that it is not \(p\)-universal in the stronger sense of
\cite[Definition 3.7]{NR2007}.  Let \(F=\ell_p^2\), realized as
\(L_p(\{1,2\})\) with counting measure, and let \(\rho\) be the trivial
representation of \(G\) on \(F\).  Then \(F\) is a \(\QSL_p\)-space, so
\((\rho,F)\in\Rep_p(G)\).  It is small: \(\card(L_1(G))=\card(\mathbb C)\),
and
\[
  \card(F)=\card(\mathbb C)\leq \card(\mathbb C)^{\aleph_0}
  = \card(L_1(G))^{\aleph_0}.
\]
Thus \((\rho,F)\) is one of the small representations that must be contained
by any \(p\)-universal representation in the sense of Definition 3.7.

But \((\rho,\ell_p^2)\) is not contained in \((\mathbf 1,\mathbb C)\), since
containment would require an isometric equivalence of \(\ell_p^2\) with a
closed subspace of the one-dimensional space \(\mathbb C\), which is
impossible by dimension.  Hence \((\mathbf 1,\mathbb C)\) is old
cyclic-universal but not small-universal.
\end{proof}

\section*{Verification Notes}

The example uses only the trivial group, so no nontrivial continuity or
integrability issue is hidden.  The representation on \(\ell_p^2\) is a
representation on an actual \(L_p\)-space, hence certainly on a \(\QSL_p\)
space.  The smallness estimate is exactly the cardinal bound in the source
definition.  The failure of containment is purely dimensional.

\section*{Novelty Check}

Local run indexes were searched for \texttt{0711.2057}, the paper title,
\(\QSL_p\), \(p\)-operator-space, and row/column keywords; no existing packet
or attempt for this source question was found.  Web searches on June 19, 2026
for phrases including ``Column and row operator spaces over QSL\_p-spaces we
do not know'', ``universal p-pseudofunctions p-universal cyclic small
representations'', and related Herz--Schur/\(A_p(G)\) terms did not locate a
published answer to this converse.  The result should be reviewed as a simple
new counterexample to the definitional converse, not as progress on the
paper's deeper multiplier questions.

\section*{Human Review Recommendation}

Check that Runde's Definition 4.5 indeed allows the one-dimensional trivial
representation to count as \(p\)-universal for \(G=\{e\}\), and that the
zero representation, if included, causes no exception.  The main conclusion
then follows immediately.

\begin{thebibliography}{9}

\bibitem{NR2007}
M. Neufang and V. Runde,
\emph{Column and row operator spaces over \(\QSL_p\)-spaces and their use in
abstract harmonic analysis},
arXiv:0711.2057.

\bibitem{Runde2004}
V. Runde,
\emph{Representations of locally compact groups on \(\QSL_p\)-spaces and a
\(p\)-analog of the Fourier--Stieltjes algebra},
arXiv:math/0402018.

\end{thebibliography}

\end{document}
